\section{The Sanctoral Cycle and the Printed Calendar}
\label{sec:sanctoral}

\subsection{How the calendar is laid out}

The universal calendar of the 1962 typical edition occupies printed pages XLV--LIII, nine pages
of five columns: the cycle of epacts, the dominical letter, the Roman day of the month by
Kalends, Nones, and Ides, the civil day of the month, and the entry itself. An entry carries the
Latin title, a qualification of the Saint's state where the Missal's abbreviations require it,
and the class. Subordinate commemorations are printed in italic beneath the entry they attach
to, or on their own line where the day has no feast; a commemoration standing alone is followed
by the abbreviation \latin{Comm.}

Two kinds of entry that are not fixed-date feasts nonetheless appear in the calendar, and they
are set as notices at the foot of the month rather than against a numeral. January carries two:
\latin{Dominica inter octavam Nativitatis Domini et Epiphaniam, vel, ea deficiente, die
2 ianuarii: Sanctissimi Nominis Iesu, II classis}, and \latin{Dominica I post Epiphaniam:
S. Familiae, Iesu, Mariae, Ioseph, II classis}. March carries the notice \latin{Feria VI post
dominicam I Passionis: Commemoratio septem Dolorum B. Mariae Virg., Comm.} October ends with
\latin{Dominica ultima octobris: D. N. Iesu Christi Regis, I classis}. These are the calendar's
own expression of \rn{17}'s exceptions and of the \latin{Variationes}' reduction of the Friday
of Passion Week to a commemoration.

February carries the leap-year note, which is a calendrical rule rather than a feast and is
easily overlooked:

\begin{studyframe}
\latin{In anno bissextili mensis februarius est dierum 29, et festum S. Matthiae celebratur die
25 februarii ac festum S. Gabrielis a Virg. Perdolente die 28 februarii, et bis dicitur Sexto
Kalendas, id est die 24 et die 25; et littera dominicalis, quae assumpta fuit in mense ianuario,
mutetur in praecedentem.}
\medskip

\textit{Working gloss.} In a leap year the month of February has 29 days, and the feast of
St Matthias is celebrated on 25 February and the feast of St Gabriel of Our Lady of Sorrows on
28 February; \latin{Sexto Kalendas} is said twice, that is, on the 24th and on the 25th; and the
dominical letter taken up in January is changed to the preceding one.
\end{studyframe}

\subsection{The I class entries of the universal calendar}

Read from the printed calendar pages, the universal calendar of 1962 contains fifteen
fixed-date entries of the I class and one Sunday-assigned one. The list is complete.

\begingroup\small
\renewcommand{\arraystretch}{1.12}
\begin{longtable}{@{}l >{\raggedright\arraybackslash}p{0.66\linewidth}@{}}
\toprule
\textbf{Date} & \textbf{Entry as printed} \\
\midrule
\endfirsthead
\toprule
\textbf{Date} & \textbf{Entry as printed} \\
\midrule
\endhead
\bottomrule
\endfoot
1 January & \latin{Octava Nativitatis Domini} \\
6 January & \latin{In Epiphania Domini} \\
19 March & \latin{S. Ioseph, Sponsi B. Mariae Virg., Conf. et Eccl. univ. Patroni} \\
25 March & \latin{Annuntiatio B. Mariae Virg.} \\
1 May & \latin{S. Ioseph Opificis, Sponsi B. Mariae Virg., Conf.} \\
24 June & \latin{In Nativitate S. Ioannis Baptistae} \\
29 June & \latin{Ss. Petri et Pauli App.} \\
1 July & \latin{Pretiosissimi Sanguinis D. N. I. C.} \\
15 August & \latin{In Assumptione B. Mariae Virg.} \\
29 September & \latin{In Dedicatione S. Michaelis Arch.} \\
1 November & \latin{Omnium Sanctorum} \\
2 November & \latin{In Commemoratione omnium Fidelium defunctorum} \\
8 December & \latin{In Conceptione Immaculata B. Mariae Virg.} \\
24 December & \latin{Vigilia} \\
25 December & \latin{In Nativitate Domini, I classis cum oct.} \\
\midrule
Last Sunday of October & \latin{D. N. Iesu Christi Regis} \\
\end{longtable}
\endgroup


To these the temporal cycle adds the movable I class celebrations that the calendar pages do not
carry: Easter and Pentecost with their octaves and their days within, Ash Wednesday, the ferias
of Holy Week, the vigil of Pentecost, the Ascension, the Most Holy Trinity, Corpus Christi, and
the Most Sacred Heart. The table at \rn{91} ranks all of them together, and Section~\ref{sec:precedence}
gives it.

\subsection{The II class entries of the universal calendar}

The II class layer is larger and structurally more varied, because it contains feasts, vigils,
and days within the octave of the Nativity in the same rank. Read from the same pages, it is:

\begingroup\small
\renewcommand{\arraystretch}{1.08}
\begin{longtable}{@{}l >{\raggedright\arraybackslash}p{0.38\linewidth} l >{\raggedright\arraybackslash}p{0.32\linewidth}@{}}
\toprule
\textbf{Date} & \textbf{Entry} & \textbf{Date} & \textbf{Entry} \\
\midrule
\endfirsthead
\toprule
\textbf{Date} & \textbf{Entry} & \textbf{Date} & \textbf{Entry} \\
\midrule
\endhead
\bottomrule
\endfoot
13 Jan & \latin{In Commemoratione Baptismatis D. N. I. C.} & 14 Aug & \latin{Vigilia} (of the Assumption) \\
2 Feb & \latin{In Purificatione B. Mariae Virg.} & 16 Aug & \latin{S. Ioachim, Patris B. Mariae Virg., Conf.} \\
22 Feb & \latin{Cathedrae S. Petri Ap.} & 22 Aug & \latin{Immaculati Cordis B. Mariae Virg.} \\
24 Feb & \latin{S. Matthiae Ap.} & 24 Aug & \latin{S. Bartholomaei Ap.} \\
25 Apr & \latin{S. Marci Evangelistae} & 8 Sep & \latin{In Nativitate B. Mariae Virg.} \\
11 May & \latin{Ss. Philippi et Iacobi App.} & 14 Sep & \latin{In Exaltatione S. Crucis} \\
31 May & \latin{B. Mariae Virg. Reginae} & 15 Sep & \latin{Septem Dolorum B. Mariae Virg.} \\
23 Jun & \latin{Vigilia} (of St John the Baptist) & 21 Sep & \latin{S. Matthaei Ap. et Evangelistae} \\
28 Jun & \latin{Vigilia} (of Ss Peter and Paul) & 7 Oct & \latin{B. Mariae Virg. a Rosario} \\
2 Jul & \latin{In Visitatione B. Mariae Virg.} & 11 Oct & \latin{Maternitatis B. Mariae Virg.} \\
25 Jul & \latin{S. Iacobi Ap.} & 18 Oct & \latin{S. Lucae Evangelistae} \\
26 Jul & \latin{S. Annae Matris B. Mariae Virg.} & 28 Oct & \latin{Ss. Simonis et Iudae App.} \\
6 Aug & \latin{In Transfiguratione D. N. I. C.} & 9 Nov & \latin{In Dedicatione Archibasilicae Ss.mi Salvatoris} \\
10 Aug & \latin{S. Laurentii Mart.} & 30 Nov & \latin{S. Andreae Ap.} \\
\midrule
21 Dec & \latin{S. Thomae Ap.} & 29 Dec & \latin{De V die infra octavam Nativitatis Domini} \\
26 Dec & \latin{S. Stephani Protomart.} (II day within the octave) & 30 Dec & \latin{De VI die infra octavam Nativitatis Domini} \\
27 Dec & \latin{S. Ioannis Ap. et Evang.} (III day within the octave) & 31 Dec & \latin{De VII die infra octavam Nativitatis Domini} \\
28 Dec & \latin{Ss. Innocentium Mart.} (IV day within the octave) & & \\
\midrule
\multicolumn{4}{@{}l@{}}{Sunday between the octave of the Nativity and the Epiphany, otherwise 2 January: \latin{Sanctissimi Nominis Iesu}} \\
\multicolumn{4}{@{}l@{}}{First Sunday after the Epiphany: \latin{S. Familiae, Iesu, Mariae, Ioseph}} \\
\end{longtable}
\endgroup


Thirty-seven II class entries stand in the calendar, of which three are vigils (23 June,
28 June, 14 August), six are days within the octave of the Nativity — three of them
simultaneously feasts — and two are assigned to Sundays. The vigil count agrees exactly with
\rn{31}, which names four II class vigils, the fourth being the movable vigil of the Ascension.
This is the kind of cross-check a reference should perform and report: the closed list in the
rubrics and the entries in the calendar account for each other without residue.

\subsection{The III class and commemorative layers}

Everything else in the universal calendar is a III class feast or a commemoration. That layer is
large — the great majority of the calendar's entries — and this edition does not publish an
entry-by-entry inventory of it. What can be said about its structure without inventorying it is
this.

III class feasts are the ordinary sanctoral grade of the 1962 system and the residue of four
older grades: \latin{duplex maius}, \latin{duplex minus}, \latin{semiduplex}, and the
post-1955 \latin{simplex} (\latin{Variationes} \rnn{3--4}). They have no First Vespers of their
own (\rn{37}\,b). In the table of precedence they stand at line 24 for universal feasts and
line 23 for particular ones, below the Lenten and Passiontide ferias at line 22 and above the
Advent ferias at line 25. They can be perpetually impeded, and \rn{100} treats them differently
according to whether the impediment is diocese-wide or affects only some churches.

Commemorations in the calendar are not feasts of a lower class. They are the mode of non-full
celebration defined at \rn{5}, and their standing entry in the calendar means that in the Office
and Mass of whatever day occurs, an oration is added, subject to the counting rules of
Section~\ref{sec:commemorations}. Missal rubrics \rn{302}\,b and \rn{303}\,b add a further
possibility: a Mass \latin{de commemoratione in Officio diei occurrente} may be said as a
festive Mass in the broader sense, but only on a IV class liturgical day.

\subsection{What 1960 changed, stated as change}

Because the 1962 Missal prints only the result, the only place where the reform's calendar
decisions are stated as decisions is the first chapter of the \latin{Variationes} annexed to the
promulgating decree. Its substantive provisions beyond the grade conversion already given are
these.

\textbf{Reduced to commemorations} (\rn{5}): St George, Martyr (23 April); Our Lady of Mount
Carmel (16 July); St Alexius, Confessor (17 July); Ss Cyriacus, Largus, and Smaragdus, Martyrs
(8 August); the Impression of the Stigmata of St Francis (17 September); Ss Eustace and
Companions, Martyrs (20 September); Our Lady of Ransom (24 September); St Thomas, Bishop and
Martyr (29 December); St Silvester I, Pope and Confessor (31 December); and the Seven Sorrows of
the Blessed Virgin Mary on the Friday after the First Sunday of the Passion.

\textbf{Raised to the I class} (\rn{6}): the octave day of the Nativity (1 January) and the
Commemoration of All the Faithful Departed (2 November), the latter with the express
qualification \latin{quae tamen locum cedere pergit dominicae occurrenti} — which nevertheless
continues to yield place to an occurring Sunday.

\textbf{Raised to the II class} (\rn{7}): the Holy Family (First Sunday after the Epiphany), the
Chair of St Peter (22 February), and the Exaltation of the Holy Cross (14 September).

\textbf{Struck from the calendar} (\rn{8}): the Chair of St Peter at Rome (18 January); the
Finding of the Holy Cross (3 May); St John before the Latin Gate (6 May); the Apparition of
St Michael the Archangel (8 May); St Leo II (3 July); St Anacletus (13 July); St Peter in Chains
(1 August); and the Finding of St Stephen (3 August). The commemoration of St Vitalis, Martyr
(28 April) is likewise struck.

\textbf{Newly inscribed} (\rn{9}): the Commemoration of the Baptism of Our Lord Jesus Christ
(13 January); St Gregory Barbarigo, Bishop and Confessor (17 June); and St Anthony Mary Claret,
Bishop and Confessor (23 October).

\textbf{Transferred} (\rnn{10--11}): St Irenaeus from 28 June to 3 July; St John Mary Vianney
from 9 to 8 August; and the commemoration of Ss Sergius, Bacchus, Marcellus, and Apuleius from 7
to 8 October.

\textbf{Renamed} (\rn{12}): the feast of the Circumcision of the Lord becomes \latin{Octava
Nativitatis Domini} (1 January); the feast of the Chair of St Peter the Apostle at Antioch
becomes simply \latin{Festum Cathedrae S. Petri Ap.} (22 February); and the feast of the Most
Holy Rosary of the Blessed Virgin Mary becomes \latin{Festum B. Mariae Virg. a Rosario}
(7 October).

The pattern is legible. Four of the eight suppressions remove a duplicate observance of a Saint
or mystery already honoured on another day — a second Chair of Peter, a second Cross, a second
Michael, a second Stephen. Two of the three renamings do the same work by other means: the
Circumcision is redescribed by its place in the octave, and the Antioch qualification is dropped
from the surviving Chair. The three new inscriptions are of a different character: two recent
canonisations and one replacement for a suppressed octave day. Nothing here is a doctrinal
judgement; it is the tidying that \latin{Rubricarum instructum} announced.

\subsection{Where a feast falls, and what happens when it cannot}

\rnn{59--62} govern the assignment of days. \rn{59}: feasts already introduced into calendars
are to be celebrated on the day on which they are now inscribed. \rn{60}, for new universal
feasts: feasts of Saints are ordinarily assigned to the \latin{dies natalicius}, the day on which
the Saint was born to eternal life; if that day is impeded from any cause, the feast is assigned
to a day determined by the Holy See, which is then held as a \latin{dies quasi-natalicius}; for
other feasts, the Holy See assigns the day.

\rn{61}, for new particular feasts, is more detailed and worth having in full because it is the
rule by which most diocesan and religious calendars are actually built:

\begin{studybox}{\rn{61}: assigning new particular feasts}
\begin{enumerate}[label=\alph*), leftmargin=1.6em, itemsep=0.2em]
\item Proper feasts of Saints or Blessed are ordinarily assigned to the \latin{dies natalicius},
unless it is impeded or the Holy See has disposed otherwise. But feasts proper to a place or
church that are also inscribed in the universal, diocesan, or religious calendar at a lower
grade are to be celebrated on the same day as in that calendar.
\item If the \latin{dies natalicius} is unknown, feasts are assigned, with the approval of the
Holy See, to a day that is of the IV class in the perpetual diocesan or religious calendar.
\item If the \latin{dies natalicius} is perpetually impeded for the whole diocese or religious
institute or proper church, then feasts of the I or II class are assigned to the nearest
following day that is not of the I or II class, and feasts of the III class to the nearest
following day free from other feasts and Offices of equal or higher grade.
\item Indult feasts are inscribed on the day assigned by the Holy See in the concession.
\end{enumerate}
\end{studybox}

\rn{62} handles the case of several Saints inscribed under one feast: they are always celebrated
together as the Breviary gives them, whenever they are to be honoured at the same grade, even if
one or some are more proper. If one or some are to be honoured at the I class, the Office is of
these alone and the companions are omitted; if one or some are more proper and to be honoured at
a higher grade, the whole Office is of the more proper with a commemoration of the companions.
